\documentclass[11pt,aspectratio=43]{beamer}

% -----------------------------------------------------------------------------
% Theme and typography
% -----------------------------------------------------------------------------
\usetheme{Madrid}
\usefonttheme{serif}

\definecolor{RecommenderBlue}{RGB}{40, 60, 140}
\setbeamercolor{structure}{fg=RecommenderBlue}
\setbeamercolor{palette primary}{bg=RecommenderBlue,fg=white}
\setbeamercolor{palette secondary}{bg=RecommenderBlue!90!black,fg=white}
\setbeamercolor{palette tertiary}{bg=RecommenderBlue!80!black,fg=white}
\setbeamercolor{palette quaternary}{bg=RecommenderBlue!70!black,fg=white}
\setbeamercolor{frametitle}{bg=RecommenderBlue,fg=white}
\setbeamercolor{title}{bg=RecommenderBlue,fg=white}
\setbeamercolor{block title}{bg=RecommenderBlue,fg=white}
\setbeamercolor{block body}{bg=black!8,fg=black}

\setbeamertemplate{navigation symbols}{}
\setbeamertemplate{itemize item}[circle]
\setbeamertemplate{itemize subitem}[triangle]
\setbeamertemplate{enumerate item}[circle]
\setbeamertemplate{caption}[numbered]
\setbeamerfont{caption}{size=\scriptsize}
\setbeamerfont{footline}{size=\tiny}

% -----------------------------------------------------------------------------
% Packages
% -----------------------------------------------------------------------------
\usepackage[T1]{fontenc}
\usepackage{amsmath,amssymb}
\usepackage{graphicx}
\usepackage{booktabs}
\usepackage{tabularx}
\usepackage{array}
\usepackage{colortbl}
\usepackage{ragged2e}

\newcolumntype{Y}{>{\RaggedRight\arraybackslash}X}
\newcommand{\strong}[1]{\textbf{#1}}

% Custom column types for tight Beamer tables
\newcolumntype{L}[1]{>{\RaggedRight\arraybackslash\hspace{0pt}}p{#1}}
\newcolumntype{Z}{>{\RaggedRight\arraybackslash\hspace{0pt}}X}

% Footer customization
\setbeamertemplate{footline}{%
  \leavevmode%
  \hbox{%
    \begin{beamercolorbox}[wd=.250\paperwidth,ht=2.25ex,dp=1ex,center]{author in head/foot}%
      \usebeamerfont{author in head/foot}\insertshortauthor
    \end{beamercolorbox}%
    \begin{beamercolorbox}[wd=.448\paperwidth,ht=2.25ex,dp=1ex,center]{title in head/foot}%
      \usebeamerfont{title in head/foot}\insertshorttitle
    \end{beamercolorbox}%
    \begin{beamercolorbox}[wd=.299\paperwidth,ht=2.25ex,dp=1ex,right]{date in head/foot}%
      \usebeamerfont{date in head/foot}\insertshortdate{}\hspace*{1.5em}%
      \insertframenumber{} / \inserttotalframenumber\hspace*{1.2ex}
    \end{beamercolorbox}%
  }%
  \vskip0pt%
}

% -----------------------------------------------------------------------------
% Metadata
% -----------------------------------------------------------------------------
\title[Bias Propagation in Recommender Loops]{Analyzing Bias Propagation in Recommender Feedback Loops Using Embeddings}
\author[Group 12]{%
  Anirudh S Nair \\
  Navya M J \\
  Niharika S \\
  Shravan Nander Pandala \\[1.5ex]
  Project Guide: Dr. Ajeesh Ramanujan%
}
\institute[College of Engineering Trivandrum]{%
  Department of Computer Science and Engineering\\
  College of Engineering Trivandrum}
\date{August 4, 2026}

\begin{document}

% Slide 1: Title
\begin{frame}
  \titlepage
\end{frame}

% Slide 2: Outline
\begin{frame}{Outline}
  \tableofcontents
\end{frame}

% -----------------------------------------------------------------------------
% Illustrative Example Slide
% -----------------------------------------------------------------------------
\begin{frame}{Feedback Loop Bias Example}
\scriptsize

\textbf{Netflix}

\[
\text{Watch Stranger Things}
\]

\vspace{-0.2cm}
\[
\Downarrow
\]

\vspace{-0.2cm}
\[
\text{Recommend more Sci-Fi}
\]

\vspace{-0.2cm}
\[
\Downarrow
\]

\vspace{-0.2cm}
\[
\text{User watches Sci-Fi}
\]

\vspace{-0.2cm}
\[
\Downarrow
\]

\vspace{-0.2cm}
\[
\text{Model learns: "Likes Sci-Fi"}
\]

\vspace{-0.2cm}
\[
\Downarrow
\]

\vspace{-0.2cm}
\[
\text{Future recommendations mostly Sci-Fi}
\]

\vspace{0.15cm}

\textbf{Feedback Loop:}
Recommendations shape user behavior, which is used to retrain the model, reinforcing the same recommendations and reducing diversity.

\end{frame}

% -----------------------------------------------------------------------------
% Section 1: Domain of Interest
% -----------------------------------------------------------------------------
\section{Domain of Interest}
\begin{frame}{Domain of Interest}
  \begin{itemize}
    \item \strong{Recommender Systems (RSs):} Digital platforms use recommendation systems to discover user preferences and match users to catalog items.
    
    \item \strong{Vector Embeddings:} Recommendation engines map users and items into a shared, low-dimensional space as dense vectors.
    
    \item \strong{Closed Feedback Loops:} Recommendation outputs shape future user interactions, which become inputs for future model training.
    
    \item \strong{Embedding Geometry:} The study focuses on how vector coordinates move and change over repeated training cycles.
  \end{itemize}
\end{frame}

% -----------------------------------------------------------------------------
% Section 2: Problem Statement
% -----------------------------------------------------------------------------
\section{Problem Statement}
\begin{frame}{Problem Statement: Feedback Loops and Bias Propagation}
  \begin{itemize}
    \item Observational training data has systemic bias and deviates from true user preferences.
    
    \item Standard models favor popular items and ignore most items in the catalog.
    
    \item Feedback loop bias is dynamic, self-reinforcing, and grows stronger across retraining cycles.
    
    \item Users with niche preferences do not receive recommendations matching their real interests.
  \end{itemize}
\end{frame}

\begin{frame}{Problem Statement: Embedding Space Drift}
  \begin{itemize}
    \item Embeddings store item attributes and user preference patterns as spatial coordinates.
    
    \item Feedback loops cause vector representations to undergo continuous geometric drift.
    
    \item Vectors move away from true user preference coordinates toward a biased focal center.
    
    \item This geometric drift contracts the embedding space, reducing catalog coverage and recommendation diversity.
  \end{itemize}
\end{frame}

% -----------------------------------------------------------------------------
% Section 3: Project Objectives
% -----------------------------------------------------------------------------
\section{Project Objectives}
\begin{frame}{Project Objectives}
  \begin{block}{Primary Objective}
    Analyze and mitigate dynamic bias propagation in closed-loop recommender systems by tracking geometric changes in vector embedding space.
  \end{block}
  
  \vspace{0.3cm}
  
  \begin{enumerate}
    \item \strong{Build Simulation Framework:} Construct a closed-loop simulation environment to model multi-generation recommendation cycles.
    
    \item \strong{Quantify Vector Contraction:} Calculate the velocity of embedding space contraction to measure diversity loss and popularity concentration.
    
    \item \strong{Evaluate Mitigation Methods:} Test alternative recommendation algorithms and bias mitigation techniques to stabilize vector trajectories.
  \end{enumerate}
\end{frame}

% -----------------------------------------------------------------------------
% Expected Outcome
% -----------------------------------------------------------------------------
\begin{frame}{Expected Outcome}
\begin{itemize}
    \item Develop a recommendation model for feedback loop analysis.
    \item Train and evaluate the model through iterative recommendation cycles.
    \item Analyze the evolution of embedding representations over time.
    \item Study the impact of feedback loops on bias and recommendation diversity.
    \item Provide insights for future bias-aware recommender systems.
\end{itemize}
\end{frame}

% -----------------------------------------------------------------------------
% Section 4: Literature Review
% -----------------------------------------------------------------------------
\section{Literature Review}

\begin{frame}{Literature Review}
\scriptsize
\setlength{\tabcolsep}{1.5pt}
\renewcommand{\arraystretch}{1.05}
\begin{tabularx}{\textwidth}{L{1.8cm} Z Z Z}
\toprule
\strong{Paper} & \strong{Methodology} & \strong{Key Findings} & \strong{Limitations} \\
\midrule
Zoralioglu \& Yalcin (2026) & 10-iteration loop framework on benchmarks. & Loops degrade calibration and long-tail item exposure. & Evaluates lists directly; lacks vector-space metrics. \\
\midrule
Liu et al. (2026) & CDSRec model using Markov chains, DAGs, and VAEs. & Separates intrinsic quality from trend bias. & High complexity; no coordinate contraction tracking. \\
\midrule
Chen et al. (2023) & Survey of debiasing (IPS, Causal, RL). & Identifies feedback loops as primary bias driver. & Qualitative; lacks empirical vector tracking. \\
\midrule
Mansoury et al. (2020) & Offline simulation across UserKNN, BPR, MostPopular. & Loops amplify popularity bias and drop diversity. & Fixed acceptance rules; no latent vector analysis. \\
\midrule
Chaney et al. (2018) & Multi-generation agent-based simulations. & Homogenizes users and harms niche preferences. & Relies on synthetic utility; ignores vector drift. \\
\bottomrule
\end{tabularx}
\end{frame}

% -----------------------------------------------------------------------------
% Section 5: Motivation
% -----------------------------------------------------------------------------
\section{Motivation}
\begin{frame}{Motivation}
  \begin{itemize}
    \item \strong{Research Gap 1: Static vs Dynamic Bias} -- Existing mitigation strategies treat bias as a static problem, ignoring how feedback loops amplify bias over time.
    
    \item \strong{Research Gap 2: Structural Vector Analysis} -- Existing studies do not measure the geometric drift of vector embeddings during retraining loops.
    
    \item \strong{Bridging the Gap:}
      \begin{itemize}
        \item Define mathematical metrics to track embedding displacement and contraction speed.
        \item Identify algorithmic choices that prevent latent space collapse.
        \item Preserve catalog coverage and improve recommendations for users with niche preferences.
      \end{itemize}
  \end{itemize}
\end{frame}

% Slide: End Slide
\begin{frame}
  \centering
  \LARGE Thank You!
\end{frame}

\end{document}